\section{Problem Statement}

Given a five-channel Cell Painting image $\mathbf{x} \in \mathbb{R}^{5 \times H \times W}$ and a
text description $t$ of the applied perturbation, we learn image and text encoders that map both
inputs to the same normalized embedding space:
\begin{equation}
  f_{\text{img}}(\mathbf{x}),\;
  f_{\text{txt}}(t) \in \mathbb{R}^d,
  \quad
  \|f_{\text{img}}(\mathbf{x})\|_2
  = \|f_{\text{txt}}(t)\|_2
  = 1.
\end{equation}
Matching image--text pairs should have high cosine similarity, while unrelated pairs should be
farther apart. We set $d=512$. Although the equation is written for one image, the model operates
at the well level by combining all imaged sites from a well into one vector.

The trained encoders support three retrieval tasks:
\begin{itemize}
  \item \textbf{Text-to-image retrieval.}
        Given a perturbation description, retrieve the wells that show the corresponding phenotype.
  \item \textbf{Image-to-text retrieval.}
        Given a well, return the perturbation description that best matches its morphology.
  \item \textbf{Image-to-image retrieval.}
        Compare well embeddings by cosine similarity, without text at inference time. This is the
        setting of the standard CPJUMP1 benchmark: do replicate wells of one perturbation match, and
        do perturbations that share a target match across modalities?
\end{itemize}

This task is harder than ordinary image--caption alignment for three reasons.

\paragraph{Batch effects.}
The same perturbation can look different across plates and imaging days. Changes in illumination,
staining or focus may be stronger than the biological effect of interest. In CPJUMP1,
reproducibility also varies by perturbation type and cell line~\cite{chandrasekaran2024jump}.

\paragraph{Imperfect labels.}
Biological similarity is not simply positive or negative. Two compounds may share a primary target
but differ in their secondary effects, while a gene knockout and a compound may influence the same
pathway through different mechanisms. This motivates softer supervision than a standard loss in
which every non-matching pair is treated as equally negative.

\paragraph{Scale and compute.}
CPJUMP1 contains several terabytes of microscopy data. Training both foundation models end to end
is impractical on a single consumer GPU, so we run the frozen backbones once, cache their outputs
and train only the smaller modules on top.
